\documentclass[letterpaper]{article}
\usepackage[spanish]{babel}
\selectlanguage{spanish}
\usepackage[utf8]{inputenc}
\usepackage{amsmath,amssymb}
\usepackage{graphicx}
\usepackage{float}
\usepackage{booktabs}
\usepackage{array}
\usepackage{hyperref}
\usepackage{enumitem}

\graphicspath{ {./figures/} }
\usepackage[left=1.25in,right=1.25in,top=1.0in,bottom=1.0in]{geometry}
\setlength{\parskip}{6pt}
\setlength{\parindent}{0pt}

\begin{document}

\begin{center}
{\Large\bf Análisis de Resultados --- Entrega 4}\\[4pt]
{\large\bf Minimización del tiempo de respuesta logística para la distribución}\\[2pt]
{\large\bf de insumos humanitarios de SENAPRED --- Región de Valparaíso}\\[4pt]
{\normalsize Grupo 92 \quad ICS1113-Optimización}
\end{center}
\hrule\vspace{6pt}
\tableofcontents
\newpage

% ──────────────────────────────────────────────────────────────────────────────
\section{Características del modelo}

El programa del modelo generó un total de \textbf{311.366 variables de decisión}. Estas se
dividen según su naturaleza matemática en tres categorías:

\begin{itemize}[leftmargin=2em, itemsep=3pt]
  \item \textbf{Variables Binarias} ($w_{it}$ e $y_{it}$): Determinan las decisiones
    estratégicas de apertura y operatividad de cada bodega candidata en cada
    período mensual. Representan el núcleo combinatorio del problema: una vez
    que $w_{it}=1$, el modelo compromete el costo fijo $F_i$ y activa la
    capacidad logística de esa instalación de forma irrevocable para el
    horizonte restante.

  \item \textbf{Variables Enteras Generales} ($n_{ijmtp}$, $e_{it}$): Codifican
    decisiones logísticas que exigen discretud física. $n_{ijmtp}$ es el número
    de viajes del vehículo $m$ entre la bodega $i$ y la comuna $j$ bajo
    prioridad $p$ en el período $t$; no se puede asignar ``medio camión''.
    $e_{it}$ es la dotación de personal, cuya integridad garantiza que la
    restricción R10a de dotación mínima opere sobre personas reales y no sobre
    fracciones.

  \item \textbf{Variables Continuas} ($x_{ijkmtp}$, $s_{ikt}$, $c_{ikt}$,
    $f_{jktp}$): Representan los flujos físicos de insumos entre nodos, los
    niveles de inventario al cierre de cada período, las compras mensuales y la
    demanda insatisfecha. Su naturaleza continua es apropiada porque los
    suministros humanitarios (unidades de agua, raciones, medicamentos) pueden
    fraccionarse a efectos de planificación sin alterar la factibilidad de la
    solución.
\end{itemize}

Toda esta red de variables está sujeta a \textbf{338.595 restricciones
matemáticas} que garantizan el balance de masa (R2--R3), el respeto a las
capacidades físicas de bodegas y vehículos (R4--R5), las condiciones de
aptitud y factibilidad temporal (R6--R8), las holguras de personal (R9--R10b),
la restricción de presupuesto (R13) y el condicionamiento de compras a
apertura (R14). La matriz de coeficientes evalúa \textbf{2.076.596 elementos
no nulos}.

Tras \textbf{1.800,39 segundos} de procesamiento (equivalente al límite
máximo de 30 minutos), el solver alcanzó un \textbf{GAP de optimalidad del
$6{,}02\%$}, reportando un valor de la función objetivo de
$Z^* = 131{,}629{,}817$ minutos-equivalentes. Un GAP del $6{,}02\%$ implica
que la solución obtenida está a lo más un $6{,}02\%$ por encima del óptimo
global desconocido, un resultado satisfactorio dada la complejidad combinatoria
del modelo con más de 311.000 variables.

% ──────────────────────────────────────────────────────────────────────────────
\section{Estrategia de localización y posicionamiento de bodegas}

\textit{Documento: 01\_Reporte\_Bodegas\_Abiertas.csv}

Al analizar los resultados numéricos obtenidos, se observa que el modelo decidió
habilitar las \textbf{12 bodegas candidatas}, agotando la totalidad de las
ubicaciones disponibles. En lugar de concentrar el presupuesto en pocas
instalaciones de gran tamaño, el programa optó por una red plenamente
descentralizada que cubre de forma simultánea cuatro zonas geográficas
estratégicas de la Región de Valparaíso:

\begin{itemize}[leftmargin=2em, itemsep=3pt]
  \item \textbf{Costa norte y sur}: Puerto Quintero y La Ligua cubren el litoral
    norte, mientras que Puerto San Antonio protege el extremo sur costero, donde
    la concentración portuaria e industrial eleva el riesgo de incendios
    de interfaz.
  \item \textbf{Polo central urbano-forestal}: Placilla (hub de la conurbación
    Valparaíso-Viña), Aeródromo Rodelillo, Belloto~Sur, Barrio Industrial
    El~Salto y Base Torquemada forman un cinturón de respuesta rápida sobre las
    zonas de cerros con mayor historial de megaincendios.
  \item \textbf{Arco interior central}: Estadio Limache y Casablanca Centro
    actúan como puntos de relay hacia comunas de interior con acceso vial
    complejo.
  \item \textbf{Valles interiores cordilleranos}: Ruta~60CH~Quillota y Parque
    Industrial San Felipe cubren el eje del Aconcagua, permitiendo atender
    comunas como Quillota, San Felipe, La~Cruz y, en forma crítica, Los~Andes,
    incorporada en esta entrega como nuevo nodo demandante.
\end{itemize}

\begin{table}[H]
\centering
\caption{Bodegas habilitadas y mes de apertura}
\label{tab:bodegas}
\begin{tabular}{lrr}
\hline
\textbf{Bodega} & \textbf{Mes apertura} & \textbf{Costo fijo (CLP)} \\
\hline
Placilla                & 1 & 59.603.887 \\
Pto\_San\_Antonio       & 1 & 45.734.667 \\
Barrio\_Ind\_El\_Salto  & 1 & 33.983.161 \\
Belloto\_Sur            & 1 & 31.854.632 \\
Base\_Torquemada        & 1 & 29.625.860 \\
Ruta\_60CH\_Quillota    & 1 & 31.744.420 \\
Parque\_Ind\_San\_Felipe & 1 & 32.360.050 \\
Pto\_Quintero           & 1 & 27.177.706 \\
Aerodromo\_Rodelillo    & 1 & 13.986.210 \\
Casablanca\_Centro      & 1 & 10.456.831 \\
Estadio\_Limache        & 1 & 12.543.632 \\
La\_Ligua               & 1 & 11.091.252 \\
\hline
\end{tabular}
\end{table}

Adicionalmente, el \textbf{100\% de estas bodegas se habilita en el mes~1}.
Esto demuestra cuantitativamente el valor de la fase de preparación
anticipada. Al comprometer el costo fijo de apertura en el primer período, la
red obtiene su máxima capacidad de almacenamiento antes de que comience la
temporada de incendios (enero-marzo), permitiendo acopiar inventario con meses
de antelación y evitando que la respuesta ante una emergencia quede condicionada
a los tiempos de compra y traslado desde proveedores externos. Esta estrategia
es además la única matemáticamente viable para satisfacer la restricción de
tiempo máximo $T^{\max}_1 = 45$~min para la prioridad~1: ninguna configuración
con menos bodegas es capaz de cubrir los 20 nodos demandantes dentro de esa
ventana temporal con los modos de transporte disponibles.

% ──────────────────────────────────────────────────────────────────────────────
\section{Estadísticas de insumos y triage}

\textit{Documento: 02\_Reporte\_Faltante.csv}

El reporte de demanda insatisfecha valida de manera contundente la efectividad
de los ponderadores de penalización de la función objetivo. Se observa una
fuerte asimetría en la atención por nivel de prioridad:

\begin{table}[H]
\centering
\caption{Demanda insatisfecha por nivel de prioridad}
\label{tab:faltante}
\begin{tabular}{lr}
\hline
\textbf{Prioridad} & \textbf{Faltante total (unidades)} \\
\hline
1 (soporte vital crítico) &     296 \\
2 (atención media)        &  10.573 \\
3 (baja urgencia)         &  97.031 \\
\hline
\end{tabular}
\end{table}

\begin{figure}[H]
\centering
\includegraphics[width=0.55\textwidth]{G1_faltante_prioridad.png}
\caption{Distribución porcentual del faltante por nivel de prioridad}
\label{fig:faltante}
\end{figure}

La prioridad~1 concentra apenas \textbf{296 unidades faltantes} en todo el
horizonte anual, mientras que las prioridades~2 y~3 absorben los quiebres de
stock con 10.573 y 97.031 unidades respectivamente. Ante las restricciones
conjuntas de capacidad vehicular, capacidad de bodega y presupuesto, el modelo
aplica un \textbf{triage logístico}: decide sacrificar cobertura de insumos de
baja urgencia para proteger el soporte vital básico de las comunas afectadas.
Insumos críticos como Agua~5L, Kits~Médicos de Trauma, Medicamentos~Crónicos y
Suplementos~Pediátricos presentan una cobertura prácticamente total en
prioridad~1.

El estrés logístico residual se concentra en las comunas donde la alta densidad
poblacional de interfaz urbano-forestal satura la capacidad de la flota durante
el verano: Viña~del~Mar, Quilpué, Villa~Alemana, Valparaíso y San~Antonio
registran los mayores déficits. Cabe destacar que este faltante no ocurre por
lejanía territorial —todas estas comunas cuentan con bodegas de apoyo a menos de
45 minutos— sino por la demanda excepcional que generan los megaincendios de
verano sobre una flota institucional de tamaño fijo. La incorporación de
\textbf{Los~Andes} como nuevo nodo demandante es particularmente significativa:
con tiempos terrestres superiores a 120~minutos desde la mayoría de las bodegas,
su cobertura de prioridad~1 depende del Helicóptero o de la bodega de
Parque~Ind~San~Felipe, situación que el modelo resuelve correctamente habilitando
ambas desde el mes~1.

% ──────────────────────────────────────────────────────────────────────────────
\section{Análisis de distribución de recursos}

\textit{Documento: 04\_Reporte\_Presupuesto.csv}

La sumatoria de los gastos monetarios realizados evidencia que el modelo
utiliza \textbf{\$4.804~MM~CLP} del presupuesto público disponible de
\$5.000~MM, equivalente al $96{,}1\%$ de utilización. La restricción
presupuestaria R13 es activa pero no completamente vinculante: la excelente
cobertura geográfica de las 12 bodegas permite satisfacer la mayor parte de
la demanda crítica sin necesidad de agotar el capital disponible.

\begin{table}[H]
\centering
\caption{Distribución del presupuesto utilizado}
\label{tab:presupuesto}
\begin{tabular}{lrr}
\hline
\textbf{Categoría} & \textbf{Gasto (CLP)} & \textbf{Porcentaje} \\
\hline
Compras     & 3.630.133.887 & 75{,}6\% \\
Sueldos     &   795.600.000 & 16{,}6\% \\
Apertura    &   340.162.308 &  7{,}1\% \\
Bodegaje    &    20.244.648 &  0{,}4\% \\
Transporte  &    17.793.550 &  0{,}4\% \\
\hline
\textbf{Total} & \textbf{4.803.934.393} & \textbf{96{,}1\%} \\
\hline
\end{tabular}
\end{table}

\begin{figure}[H]
\centering
\includegraphics[width=0.6\textwidth]{G3_presupuesto.png}
\caption{Distribución porcentual del presupuesto por categoría}
\label{fig:presupuesto}
\end{figure}

La adquisición de inventario (\textbf{Compras, 75{,}6\%}) es el ítem
dominante, coherente con la necesidad de pre-posicionar suministros en las 12
bodegas antes de la temporada crítica. Para lograr los tiempos de respuesta
exigidos, el modelo prioriza la disponibilidad física de insumos sobre la
eficiencia financiera: es más costoso llegar tarde que comprar antes.

Un elemento notable es el peso de los \textbf{Sueldos (16{,}6\%, \$796~MM)}.
Al habilitar las 12 bodegas desde el mes~1, la restricción R10a obliga a
mantener la dotación mínima de personal en cada instalación activa durante los
12 períodos. Esta inversión refleja que el modelo decide garantizar la
\textbf{disponibilidad operativa 24/7} de las cuadrillas: cuando se declara una
alerta, los camiones deben poder ser cargados de inmediato, sin esperar la
contratación de personal de emergencia.

El \textbf{costo de transporte (0{,}4\%, \$17{,}8~MM)} es prácticamente
despreciable. Esto es consecuencia directa de la estrategia descentralizada:
al posicionar bodegas a menos de 45~minutos de cada comuna, la distancia
promedio por viaje es mínima, lo que reduce el costo variable por ruta. La
red paga el costo de la proximidad en infraestructura y personal, pero se
ahorra sustancialmente en operación de flota.

% ──────────────────────────────────────────────────────────────────────────────
\section{Capacidad de cuadrillas y saturación}

\textit{Documento: 05\_Reporte\_Personal.csv}

El análisis de la saturación del personal revela un comportamiento
contraintuitivo pero lógico para la naturaleza del problema humanitario. Con 12
bodegas distribuidas estratégicamente en el territorio, cada instalación cubre
un área geográfica reducida y sirve a un subconjunto acotado de comunas, lo que
implica que el volumen de despachos por bodega es significativamente menor al
que concentraría una red centralizada.

La tasa promedio de saturación es \textbf{inferior al 2\%} en todas las
instalaciones, con peak en bodegas de alta actividad como Parque~Ind~San~Felipe
y Ruta~60CH~Quillota, que atienden los valles interiores con mayor diversidad
de comunas. Esto demuestra que la red opera con una \textbf{holgura logística
casi total}, asegurando que los vehículos puedan ser cargados inmediatamente al
declararse una alerta, sin generar cuellos de botella en bodega ni tiempos de
espera en el despacho.

Este patrón es una consecuencia deliberada del modelo: al invertir en
dotaciones mínimas garantizadas por R10a en las 12 instalaciones, el
optimizador sacrifica eficiencia de recursos humanos en pos de
\textbf{velocidad de respuesta}. Una bodega con personal al 2\% de saturación
no es ineficiente; es una bodega lista para operar al 100\% de su capacidad
de despacho en los primeros 45 minutos de una emergencia.

% ──────────────────────────────────────────────────────────────────────────────
\section{Despliegue de vehículos y eficiencia}

\textit{Documento: 06\_Reporte\_Rutas.csv}

El ruteo de vehículos no se realiza al azar, sino que el programa selecciona el
tipo de transporte que maximiza la eficiencia temporal para cada combinación de
ruta y prioridad, según los conjuntos de factibilidad $\mathcal{F}_p$.

\begin{table}[H]
\centering
\caption{Top 10 rutas por tasa de llenado vehicular}
\label{tab:rutas}
\begin{tabular}{llllrr}
\hline
\textbf{Origen} & \textbf{Destino} & \textbf{Vehículo} & \textbf{Mes} & \textbf{Viajes} & \textbf{Fill Rate (\%)} \\
\hline
Ruta\_60CH\_Quillota    & Quillota    & Camioneta\_4x4 &  2 & 1 & 100{,}0 \\
Ruta\_60CH\_Quillota    & San\_Felipe & Camioneta\_4x4 &  3 & 1 & 100{,}0 \\
Parque\_Ind\_San\_Felipe & San\_Felipe & Camioneta\_4x4 &  3 & 1 & 100{,}0 \\
Parque\_Ind\_San\_Felipe & San\_Felipe & Camioneta\_4x4 & 12 & 2 & 100{,}0 \\
Parque\_Ind\_San\_Felipe & Los\_Andes & Camioneta\_4x4 &  1 & 3 & 100{,}0 \\
Parque\_Ind\_San\_Felipe & Los\_Andes & Camioneta\_4x4 &  1 & 1 & 100{,}0 \\
Parque\_Ind\_San\_Felipe & Los\_Andes & Camioneta\_4x4 &  2 & 2 & 100{,}0 \\
Parque\_Ind\_San\_Felipe & Los\_Andes & Camioneta\_4x4 &  3 & 3 & 100{,}0 \\
Ruta\_60CH\_Quillota    & La\_Cruz   & Camioneta\_4x4 &  2 & 1 & 100{,}0 \\
Ruta\_60CH\_Quillota    & La\_Cruz   & Camioneta\_4x4 &  3 & 1 & 100{,}0 \\
\hline
\end{tabular}
\end{table}

\begin{figure}[H]
\centering
\includegraphics[width=0.75\textwidth]{G2_fill_rate_vehiculo.png}
\caption{Tasa de uso de capacidad por tipo de vehículo (fill rate \%)}
\label{fig:vehiculos}
\end{figure}

La \textbf{Camioneta~4x4} es el vehículo con el desempeño más consistente y
elevado: su mediana de fill rate se sitúa cerca del $57\%$, con un tercer
cuartil cercano al $93\%$, lo que significa que en la mitad superior de sus
viajes opera por encima del umbral óptimo del 80\%. Su baja capacidad
volumétrica ($2{,}5~\text{m}^3$) lejos de ser una limitación, es su ventaja
operativa: se carga rápido, transita con agilidad por las rutas de cerros y
llega dentro de los $45$~minutos, lo que la convierte en el recurso dominante
para despachos de \textbf{prioridad~1}.

El \textbf{Helicóptero} exhibe el patrón más peculiar: su rango intercuartílico
abarca desde el $13\%$ hasta aproximadamente el $96\%$, con mediana en torno
al $20\%$. Esto refleja que el modelo lo activa \emph{selectivamente} en las
rutas donde los modos terrestres no pueden cumplir $T^{\max}_1 = 45$~min
(Los~Andes, Llaillay). Cuando se activa, lo hace con alta ocupación ---el tercer
cuartil roza el $96\%$ y el máximo alcanza el $100\%$--- lo que demuestra que
el algoritmo no lo desperdicia: si el helicóptero vuela, va cargado.

El \textbf{Camión~Pesado} presenta el comportamiento más extremo: su mediana
es apenas del $5\%$, lo que indica que la mayoría de sus despachos son de
volumen reducido, pero su tercer cuartil sube al $70\%$ y llega al $100\%$ en
sus picos de demanda de verano. Es el vehículo de alto volumen reservado para
los meses críticos de enero-marzo, permaneciendo subutilizado el resto del año.
El \textbf{Camión~3/4} comparte este carácter bimodal aunque de forma menos
extrema: mediana en torno al $22\%$ y cola superior que alcanza el $100\%$.
Ambos camiones reflejan la estacionalidad del riesgo: infrautilizados en
invierno, operando a plena capacidad cuando los megaincendios de verano
generan picos de demanda que la Camioneta~4x4 sola no puede absorber.

% ──────────────────────────────────────────────────────────────────────────────
\section{Inventario y balance temporal}

\textit{Documento: 03\_Reporte\_Inventario.csv}

\begin{figure}[H]
\centering
\includegraphics[width=0.85\textwidth]{G4_stock_vs_faltante.png}
\caption{Stock final vs.\ demanda insatisfecha por mes}
\label{fig:stock}
\end{figure}

El gráfico revela con claridad la \textbf{estrategia de acumulación anticipada}
del modelo frente a la estacionalidad del riesgo de incendios.

La \textbf{demanda insatisfecha} (línea roja, eje derecho) alcanza sus máximos
en los meses de \textbf{verano}: pica en enero (mes~1) con aproximadamente
25.000 unidades faltantes, sube aún más en febrero (mes~2) con cerca de
30.000 unidades —el peor mes del horizonte—, y se mantiene elevada en
marzo (mes~3) con alrededor de 22.000 unidades. A partir de abril (mes~4)
el faltante cae abruptamente a menos de 1.000 unidades y permanece
prácticamente en cero durante los meses de invierno (mayo--octubre).
En diciembre (mes~12) el faltante vuelve a repuntar, anticipando el inicio
de la nueva temporada de incendios.

El \textbf{stock de no perecibles} (barras teal, eje izquierdo) muestra un
comportamiento opuesto y complementario. Durante enero--marzo el inventario
se consume rápidamente para atender la demanda de emergencia, manteniéndose
en niveles bajos ($\approx 150$--$300$ unidades). Entre los meses de invierno
(mayo--septiembre) el stock se reconstituye gradualmente a medida que las
bodegas realizan compras preventivas con la holgura presupuestaria disponible.
El pico de acumulación se produce en \textbf{noviembre (mes~11)}, donde el
stock total de no perecibles supera las 2.000 unidades: es el punto de máxima
preparación antes de la temporada crítica, y refleja la decisión explícita del
modelo de pre-posicionar inventario con un mes de antelación.

Los \textbf{insumos perecibles} (Agua~5L, Raciones, Medicamentos, Suplementos)
muestran stock cero al cierre de cada período en todas las bodegas, en
cumplimiento estricto de R3: el modelo los compra y despacha íntegramente
dentro del mismo mes, garantizando calidad y evitando pérdidas por vencimiento.

La concentración del faltante en verano valida que el modelo no falla por mala
planificación ni por déficit de bodegas: las 12 instalaciones y sus dotaciones
están activas desde el mes~1. El problema es la magnitud objetiva del fenómeno:
los megaincendios de enero--marzo generan una demanda de insumos que supera
la capacidad de la flota y el presupuesto disponible, independientemente de la
eficiencia logística del sistema.

\end{document}
